%=============================================================================!
% 8.A  ANSWERS TO THE CHECKS                                                  !
%=============================================================================!
\unnumbered{Answers to the checks}
\label{sec:pe-answers}

\answerto{sec:pe-origin}The size grows by a factor of 4, so the time by
$4^3 = 64$: about \SI{6.4}{\second}.

\answerto{sec:pe-classical}No. Its quantum step is $\hbar\omega =
2\pi\hbar c\,\tilde\nu = 1.2398\times10^{-4}\times3000 = \SI{0.372}{\electronvolt}$,
14.4 times $\kB T = \SI{0.02585}{\electronvolt}$.

\answerto{sec:pe-pairs}$\tau = \sigma\sqrt{m/\varepsilon}$. Krypton's
$\sigma$ is larger by a factor $3.6/3.4 = 1.06$ and its $m/\varepsilon$ by
$(83.8/39.948)/(160/120) = 1.57$, whose square root is 1.25, so its
$\tau$ is about 1.33 times argon's, longer: it is heavier, and the extra
depth of its well does not make up for that.

\answerto{sec:pe-forces}At $r = \sigma$, closer than the minimum at
$1.12\sigma$, the pair repels, $\varphi'(\sigma) = 4\varepsilon(-12 + 6)/
\sigma = -24\varepsilon/\sigma < 0$, so $\mathcal{W} = -\sigma\varphi'(\sigma) =
+24\varepsilon$: positive.

\answerto{sec:pe-cutoffs}$\lvert\varphi(2.5\sigma)\rvert = 0.0163\times
\SI{0.01034}{\electronvolt} = \SI{1.69e-4}{\electronvolt}$ per crossing.

\answerto{sec:pe-manybody}Each bond is worth $\xi/\sqrt{z}$, so the ratio
is $\sqrt{12}/\sqrt2 = \sqrt6 = 2.45$.

\answerto{sec:pe-bonded}$1 + \cos\psi$ is least, 0, at $\psi = \ang{180}$;
$1 + \cos3\psi$ is 0 where $\cos3\psi = -1$, that is $3\psi = \ang{180}$,
\ang{540} or $-\ang{180}$, so $\psi = \ang{60}$, \ang{180} or $-\ang{60}$.
Both are least at \ang{180}, the trans arrangement.

\answerto{sec:pe-coulomb}The sites $(a, b, c)$ with $a^2 + b^2 + c^2 = 2$
have two of $a$, $b$, $c$ equal to $\pm1$ and one 0: three choices of the
zero and four of the signs, 12 ions. Their $a + b + c$ is even, so they
carry the same charge as the ion at the origin.
